\section{Introduction}

LLM-powered browser agents have emerged as a promising approach to automating complex web-based workflows. Systems such as Browser-Use~\cite{browseruse}, SeeAct~\cite{seeact}, WebRL~\cite{webrl}, and WebAgent~\cite{webagent} demonstrate that language models can navigate web pages, fill forms, extract data, and complete multi-step tasks with varying degrees of success. However, three fundamental limitations persist across existing systems:

\begin{enumerate}[leftmargin=*,itemsep=2pt]
    \item \textbf{Brittle execution.} When web pages change their DOM structure, CSS selectors, or interaction patterns, pre-recorded automation scripts or skills break. Current systems require manual re-authoring or restart from scratch.
    \item \textbf{No cross-session learning.} Each task execution is effectively stateless. Successful strategies discovered in one session are not retained, catalogued, or reused in future sessions, leading to redundant LLM calls and repeated failures.
    \item \textbf{Single-agent limitations.} Complex multi-phase tasks (e.g., ``research a competitor, extract pricing, update our database, and schedule a daily monitor'') exceed the effective context window of a single agent and benefit from parallel decomposition.
\end{enumerate}

OpenSkynet is a modular, extensible framework designed to address each of these limitations. Its key contributions are:

\begin{itemize}[leftmargin=*,itemsep=2pt]
    \item \textbf{A self-healing skill executor} that captures browser state (screenshot + DOM snapshot) at failure points, sends this multimodal context to an LLM for root cause diagnosis, and automatically patches the skill with corrected steps. Each healing cycle creates a versioned update, forming a continuous improvement loop without human intervention.
    \item \textbf{An automated skill learning pipeline} that observes agent trajectories, extracts reusable skills via LLM-based review with iterative refinement (up to 3 cycles), loads recent failure trajectories as negative examples, and merges similar skills to prevent duplication. Skills can also be learned from screen recordings via multimodal trace-to-skill conversion.
    \item \textbf{A three-tier memory architecture} with working, session, and long-term memory tiers. Entries are automatically promoted based on heuristic importance scoring and LLM-driven background review. Retrieval uses hybrid vector search (OpenAI $\rightarrow$ FastEmbed $\rightarrow$ TF-IDF fallback) combined with keyword overlap scoring.
    \item \textbf{A multi-agent subagent system} with 7 built-in specialized agents (browser, code, debug, explore, integrate, review, redteam) supporting permission-based tool isolation, browser sandboxing, parallel execution (up to 3 concurrent), and explicit nesting depth control to prevent infinite recursion.
    \item \textbf{The Evil Skill Generalization Benchmark}, a deliberately adversarial test suite of 60 tasks across 20 web application templates, each featuring 7 layers of anti-LLM obfuscation designed to stress-test skill generalization and self-healing capabilities.
\end{itemize}

The remainder of this paper is organized as follows. Section~\ref{sec:architecture} describes the overall system architecture and agent loop. Section~\ref{sec:skills} details the skill learning and self-healing subsystems. Section~\ref{sec:memory} presents the tiered memory architecture. Section~\ref{sec:subagents} covers the multi-agent subagent system. Section~\ref{sec:benchmark} introduces the Evil Skill Generalization Benchmark. Section~\ref{sec:related} discusses related work, and Section~\ref{sec:conclusion} concludes with directions for future work.
